\chapter{Точная LCF-проверка родительского действия Краус-моста}
\label{ch:v8-kraus-bridge-parent-action-lcf-migration-gate}

\section{Полевая форма без численного допуска}

Для нормированного центрального контраста и каждого из двенадцати
межсекторных направлений eDSL точно вычисляет
\begin{equation}
 -\langle Q,\mathcal L_aQ\rangle=\frac7{36}.
\end{equation}
Восемь внутренних лептонных контролей дают ноль. Поэтому
\begin{equation}
 E_{q\ell}(z)=\frac7{36}\sum_{a=1}^{12}z_a^2,
 \qquad
 \Hess E_{q\ell}=\frac7{18}I_{12}.
 \label{eq:v8-lcf-kraus-parent-hessian}
\end{equation}
После вложения в использованный исходным гейтом 27-мерный вещественный срез
сигнатура этой добавки равна $(0,15,12)$.

\section{Символическая совместимость с двумя стационарными точками}

Прежний аудит проверял шесть значений веса. Теперь введён символ
$\lambda_{\rm bridge}\ge0$. На исходной точке точная диагональная форма
содержит семь отрицательных корневых направлений и двадцать положительных
тяжёлых направлений. На целевом вакууме положительны все 27 направлений.
Добавка \eqref{eq:v8-lcf-kraus-parent-hessian} действует только на двенадцать
уже положительных тяжёлых координат. Ядро выдаёт для всего диапазона
\begin{equation}
 \operatorname{sign}H_0(\lambda)=(7,0,20),
 \qquad
 \operatorname{sign}H_*(\lambda)=(0,0,27),
 \qquad \lambda\ge0.
 \label{eq:v8-lcf-kraus-parent-signatures}
\end{equation}

\section{Точный древесный запрет}

В вакууме $z_0=\cdots=z_{11}=0$. Поэтому одновременно
\begin{equation}
 E_{q\ell}(0)=0,
 \qquad
 \nabla E_{q\ell}(0)=0,
 \qquad
 (z_0^2,\ldots,z_{11}^2)=0.
\end{equation}
Последнее равенство непосредственно зануляет все коэффициенты
полезависимого GKSL-генератора. Положительный гессиан означает стоимость
флуктуации, но не её ненулевую интенсивность.

Если отдельно задать положительные ковариации двух семейств, точная
центральная скорость равна
\begin{equation}
 \Gamma_{q\ell}=\frac76(c_Q+c_X).
\end{equation}
Гауссовский пробный рецепт на тяжёлом гессиане при единичной внешней силе
даёт $35/96$, однако общий множитель силы этим рецептом не определяется.

\begin{theorem}[Совместимый родительское слагаемое без древесного запуска]
На точном 27-мерном срезе положительная форма Краус-моста сохраняет
сигнатуры Тома VII при каждом неотрицательном весе. В классическом целевом
вакууме все полезависимые Краус-веса равны нулю, поэтому сам межсекторный
процесс на древесном уровне выключен.
\end{theorem}

\begin{nogocriterion}[Масса не является скоростью шума]
Нельзя превращать коэффициент $7/18$ гессиана или пробное число $35/96$ в
физическую скорость. Для ненулевого канала по-прежнему требуется выведенная
ковариация, состояние среды или динамическая дилатация.
\end{nogocriterion}

\begin{protocol}
\begin{enumerate}
 \item коэффициент одного межсекторного направления: \textbf{$7/36$};
 \item гессиан межсекторного поля: \textbf{$7I_{12}/18$};
 \item сигнатура добавки в 27 измерениях: \textbf{$(0,15,12)$};
 \item диапазон веса: \textbf{все $\lambda\ge0$};
 \item исходная сигнатура: \textbf{$(7,0,20)$};
 \item вакуумная сигнатура: \textbf{$(0,0,27)$};
 \item древесные Краус-веса: \textbf{точно ноль};
 \item ковариационная скорость: \textbf{$7(c_Q+c_X)/6$};
 \item пробное значение при единичной силе: \textbf{$35/96$};
 \item физическая сила флуктуаций: \textbf{не выведена}.
\end{enumerate}
\end{protocol}

Машинный сертификат сохранён в
\path{s2t_v8_kraus_bridge_parent_action_lcf_migration_gate_results.json}.